\documentclass[12pt]{article}
\usepackage{comment,amsmath, amssymb, graphicx, caption, subcaption, float, hyperref}

\title{Modernizing Herschel's Star Gauging:\\
A Computational Approach Using Lick Observatory Data}
\author{Colin Lambert}
\date{\today}

\begin{document}
\maketitle

\begin{abstract}
This project revisits and modernizes William Herschel's star gauging method using digital imaging and contemporary computational tools. Multiple automated star-counting approaches are developed, tested, and compared using images from Lick Observatory and related sources. The goal is to evaluate detection accuracy, identify methodological limitations, and assess how modern techniques refine star-count--based stellar density estimation.
\end{abstract}

\section{Introduction}

\subsection{Context}
William Herschel’s star gauging method tried to infer the large-scale structure of the Milky Way by counting stars in different locations in the sky. His approach relied on assumptions of uniform stellar distribution and constant detectability across the telescope's field of view. While these assumptions are now known to be inaccurate, the underlying concept remains a useful framework for understanding galactic structure.

This project recreates and updates Herschel's methodology using modern digital images, telescopes, and automated star detection algorithms. By comparing multiple computational approaches, this work evaluates the feasibility and accuracy of star-count--based density estimation with contemporary instruments.

\subsection{Contributions}
This project includes:
\begin{itemize}
    \item Development of automated star-counting tools for Lick Observatory data
    \item Comparison of multiple computational detection methods
    \item Estimation of relative stellar density from detection counts
    \item Evaluation of accuracy, consistency, and methodological limitations
\end{itemize}

\subsection{Goals}
The report addresses the following questions:
\begin{itemize}
    \item Which computational methods perform best for automated star counting?
    \item How consistent are these methods under realistic observational constraints?
    \item What challenges arise from sensor noise, resolution limits, and observing conditions?
\end{itemize}

\section{Background}

\subsection{Star Gauging}
Star gauging involves counting stars within defined fields of view to infer spatial variations in stellar density. Herschel's manual counts led to one of the earliest conceptual models of the Milky Way. Despite its limitations, the method represents an important historical foundation of Galactic astronomy.

\subsection{Modern Constraints and Data Challenges}
Digital imaging introduces several challenges for automated detection:
\begin{itemize}
    \item \textbf{Limiting magnitude:} faint sources fall below detection thresholds
    \item \textbf{Noise:} read noise, dark current, and sky background
    \item \textbf{Image resolution:} atmospheric seeing and undersampling
    \item \textbf{Sensor artifacts:} hot pixels, column defects, amplifier glow
    \item \textbf{Cosmic rays:} single-frame bright artifacts
    \item \textbf{Variability:} exposure differences and tracking errors
\end{itemize}
These all arise from the use of the types of telescopes (CCDs), and are difficult to fix, especially with issues on the Lick telescope along with other things on the Unistellars.

\subsection{Computational Approaches}
Common star detection techniques include:
\begin{itemize}
    \item Threshold-based peak detection
    \item Aperture photometry and centroiding
    \item PSF-fitting and model-based extraction
    \item Segmentation algorithms (e.g., DAOStarFinder)
    \item Machine learning approaches, including CNN-based classifiers
\end{itemize}

\section{Data Sources and Preprocessing}

\subsection{Image Sources}
The dataset consists of images from Lick Observatory and potentially Unistellar telescopes, spanning a range of fields, exposure times, and observing conditions. Calibration frames may accompany some datasets.

\subsection{Raw Data Issues}
Common issues in our data include:
\begin{itemize}
    \item Noise and hot pixels
    \item Saturation and blooming
    \item Cosmic ray strikes
    \item Minor tracking errors and star elongation
\end{itemize}

\subsection{Preprocessing}
The following preprocessing steps were applied prior to detection:
\begin{itemize}
    \item \textbf{Background Estimation:} Sigma-clipped statistics to estimate background mean, median, and noise
    \item \textbf{Background Subtraction:} Median background removal to improve contrast
    \item \textbf{Data Conversion:} FITS data converted to 32-bit floating point
    \item \textbf{Normalization (CNN only):} Pixel values scaled to [0, 1]
    \item \textbf{Dimensional Handling:} Automatic selection of the first 2D slice from multi-dimensional FITS files
    \item \textbf{Noise Estimation:} Both sigma-clipped standard deviation and MAD for robustness
\end{itemize}

Traditional calibration steps (bias, dark, flat correction, cosmic ray removal) were assumed to have been applied during earlier reduction stages. This project focuses on detection and classification rather than raw calibration.

\section{Methods}

\subsection{Baseline Detection Methods}

\textbf{Thresholding with Connected Components:}
Percentile-based thresholding (95th percentile) combined with connected-component labeling was used to identify star-like regions. This approach:
\begin{itemize}
    \item Applied binary masking and region labeling
    \item Filtered detections by minimum area
    \item Found centroids via center-of-mass calculations
    \item Visualized detections using contour overlays
\end{itemize}

\textbf{Photutils Parameter Exploration:}
\texttt{DAOStarFinder} was used to explore detection behavior across parameter ranges:
\begin{itemize}
    \item Varied detection thresholds and FWHM values
    \item Used sigma-clipped background statistics
    \item Visualized detections with circular markers
    \item Logged results to text files
\end{itemize}

\subsection{Photutils--GMM Pipeline}
A more advanced pipeline combined traditional detection with machine learning:
\begin{itemize}
    \item \textbf{Detection:} \texttt{DAOStarFinder} with sigma-clipped background estimation
    \item \textbf{Feature Extraction:} Flux, peak intensity, second moments, and eccentricity
    \item \textbf{Classification:} Gaussian Mixture Model (GMM) applied to standardized features
    \item \textbf{Spatial Merging:} DBSCAN and single-link clustering to merge nearby detections
    \item \textbf{Output:} Diagnostic plots and CSV summaries with full detection metrics
\end{itemize}

DBSCAN was not that effective at fixing duplicate detections in crowded regions or in especially large objects, this will need to be improved later.

\section{Conclusion}

\subsection{Overall Assessment}
Among the tested methods, the Photutils--GMM pipeline provided the best combination of detection accuracy, robustness, and efficiency. By combining existing tools with machine learning classification, this method improved upon simpler techniques.

Key strengths include:
\begin{itemize}
    \item Reliable detection in moderately crowded fields
    \item Improved discrimination between stars and noise artifacts
    \item Efficient performance relative to connected-component methods
    \item Detailed diagnostic outputs for further analysis
\end{itemize}

\subsection{Relation to Herschel's Method}
The progression from manual counting to automated, feature-based pipelines reflects the methodological evolution since Herschel’s era. While modern tools address limitations in subjectivity and scalability, the core objective remains unchanged: using star counts to infer galactic structure.

\subsection{Limitations}
Several challenges persist:
\begin{itemize}
    \item Bias toward brighter stars
    \item Difficulty resolving severe crowding
    \item Sensitivity to parameter tuning and data quality
\end{itemize}

\subsection{Future Work}
Potential improvements include:
\begin{itemize}
    \item Enhanced preprocessing and noise modeling
    \item Improved blending models for crowded fields
    \item Completion and evaluation of CNN-based detection
    \item Calibration against external catalogs such as Gaia or Astrometry
    \item Application to multi-instrument and multi-wavelength datasets
\end{itemize}

Overall, this project demonstrates that modern computational tools can meaningfully extend Herschel’s original star gauging concept, enabling systematic and reproducible stellar density estimation with contemporary data.

\subsection{Extra Information}

All code and extra images can be viewed at my Github repository (updates are still currently in progress).
\url{https://github.com/CoolBro4811/ASTRAL_Project}

\end{document}
